% ===========================================================================
%  sviluppatore/05-estendere/06-trappole.tex
% ===========================================================================
\section{«L'ho scritto e non compare»}\label{sez:trappole}

Questa sezione è un elenco di cause, ordinate per frequenza. Serve quando un
artefatto è stato scritto, il file è al suo posto, e nell'interfaccia non
c'è.

\begin{center}
\begin{tabular}{@{}L{54mm}L{84mm}@{}}
\toprule
\textbf{Causa} & \textbf{Come si riconosce} \\
\midrule
\rigaTenue
Il registro è vecchio
  & Il processo era già avviato. Il registro di forma~B si costruisce
    all'import: chiamare \cod{reload\_registry()}, o riavviare. \\
Il nome della cartella comincia con \cod{\_}
  & È l'interruttore di spegnimento. Si vede guardando, ma è anche la cosa
    che non si guarda. \\
\rigaTenue
Manca il file atteso
  & \file{strategy.py}, \file{approach.py}, \file{weight\_mode.py},
    \file{overlay.py}. Il filtro delle cartelle lo richiede per nome. \\
Manca \cod{META} o \cod{compute}
  & La condizione di ammissione al registro li richiede entrambi. Senza, il
    modulo viene scartato senza una riga di registro. \\
\rigaTenue
\cod{META} non è fatto di soli letterali
  & Vale per i regimi, letti con \cod{ast.literal\_eval}: un \cod{np.pi}
    dentro i \cod{defaults} rende la regola invisibile. \\
La \cod{key} non coincide col nome del file
  & Rifiutata per i regimi; altrove l'artefatto compare con un nome diverso
    da quello che si sta cercando. \\
\rigaTenue
In \pkgOverlay{} si è esposto \cod{APPROACH\_DESCRIPTOR}
  & Il registro degli overlay cerca \cod{META}+\cod{compute}. Quel file è un
    \emph{approccio}, e come overlay non esiste. \\
\cod{stream\_type} fuori allowlist
  & Per gli approcci è l'unico caso che dà un \cod{ValueError} chiaro al
    caricamento. \\
\rigaTenue
\cod{output\_path\_pattern} vuoto
  & L'artefatto c'è, ma lo \emph{stream} non compare in \voce{Streams}: la
    dashboard salta i descrittori senza pattern. \\
Un altro modulo \file{ui\_ext\_*} è rotto
  & Se l'artefatto si vede ma manca il pannello che lo mostra, è questo
    (\vediSez{ui-ext}). \\
\bottomrule
\end{tabular}
\end{center}

\subsection{Il nodo della ricarica}

Il registro di forma~B si costruisce una volta, all'import del package. Un
artefatto scritto a runtime --- dallo Studio, dall'agente, dal marketplace ---
esiste su disco e non compare. \cod{reload\_registry()} risolve, ma deve fare
due cose che non sono ovvie:

\begin{codiceInterno}{python}{Le due accortezze}{overlays/\_\_init\_\_.py:144}
    Due accortezze necessarie perche' la ricarica sia davvero visibile:
    - COMPUTERS viene aggiornato SUL POSTO: `reporting.ensemble` ne tiene un
      riferimento preso all'import, e rilegarlo qui non lo aggiornerebbe;
    - EnsembleEngine.OVERLAY_REGISTRY e' un attributo di classe legato alla
      creazione della classe: va riassegnato esplicitamente.
\end{codiceInterno}

\noindent
Per le modalità di pesatura c'è un terzo punto, ed è quello che produce il
sintomo più sconcertante:

\begin{codiceInterno}{python}{Il terzo oggetto da aggiornare}{weights/\_\_init\_\_.py:117}
    try:                      # l'attributo di classe e' legato alla creazione
        from reporting.ensemble import EnsembleEngine
        EnsembleEngine.WEIGHTING_MODE_REGISTRY = REGISTRY
        # WEIGHTING_MODES e' una tupla DERIVATA dal registro alla creazione
        # della classe ed e' quella usata per validare le richieste: senza
        # aggiornarla, un weight mode nuovo compare nell'elenco ma il compute
        # risponde 400 "Invalid weighting_mode".
        EnsembleEngine.WEIGHTING_MODES = tuple(m["key"] for m in REGISTRY)
\end{codiceInterno}

\begin{attenzione}{Compare nell'elenco e risponde \cod{400}}
È il sintomo di una ricarica parziale: l'elenco che l'interfaccia mostra
viene da \cod{WEIGHTING\_MODE\_REGISTRY}, la validazione della richiesta da
\cod{WEIGHTING\_MODES}. Aggiornare il primo e non il secondo produce una voce
cliccabile che risponde «modalità non valida». Chi aggiunge un registro
nuovo di forma~B deve chiedersi quali altri oggetti ne derivino una copia al
momento della creazione della classe.
\end{attenzione}

\subsection{Due package che sono solo rinvii}

\file{reporting/weighting/} e \file{reporting/gates/} non contengono logica:
ri-esportano \pkgPesi{} e \pkgOverlay. Sono il ricordo di dove quelle due
famiglie stavano prima. Non vanno estesi, e il codice nuovo importa dalle
cartelle vere.

\subsection{Un modulo morto da non seguire}

Accanto al package \file{regime/logics/} esiste un file
\file{regime/logics.py}, \num{237} righe, con una bella tabella della
semantica delle etichette nel docstring. \textbf{Non è raggiungibile}: a
parità di nome il package vince sempre, e \cod{regime.logics} risolve a
\file{regime/logics/\_\_init\_\_.py}. Chi lo modifica non cambia niente e non
se ne accorge.

\subsection{Due incoerenze nelle mappe dell'interfaccia}

\begin{itemize}
  \item \file{dashboard/studio\_browser.py} dichiara \cod{overlay} e
        \cod{weightmode} come artefatti «a file», mentre tutti e ventuno
        quelli installati sono a cartella.
  \item \file{dashboard/agent.py} ha la stessa incoerenza al suo interno, e i
        suoi commenti la annotano già: «\emph{Qui indicava un file singolo, e
        altrove \cod{approach.py}: tre posti, tre risposte diverse alla stessa
        domanda}».
\end{itemize}

\noindent
Non rompono niente oggi, perché la scoperta accetta entrambe le forme. Vanno
sapute quando si aggiunge una famiglia: la mappa canonica da seguire è
\cod{\_FILE\_ZONA} in \file{dashboard/agent.py}.

\begin{codiceInterno}{python}{Il file principale atteso per famiglia}{dashboard/agent.py:66}
_FILE_ZONA = {"strategies": "strategy.py", "approaches": "approach.py",
              "overlays": "overlay.py", "weights": "weight_mode.py",
              "metrics": "metric.py", "indicators": "indicator.py"}
\end{codiceInterno}

\subsection{Un elenco obsoleto}

\file{strategies/\_\_init\_\_.py} elenca nel docstring «momentum, quality\_lv,
value, custom\_example». Delle quattro, due sono oggi spente (\cod{\_value},
\cod{\_custom\_example}) e mancano le altre sette attive. È documentazione
nel codice che ha smesso di essere vera, e il modo in cui ci si accorge è
cercarvi una strategia che esiste e non trovarla.
